%% GERADO por gera_tab_validacao.py. Nao editar a mao.
\begin{table}[p]
\centering
\caption{\textbf{What each source work reports as evidence for its own visual criterion.} Read from the full text of each paper; the excerpt is verbatim. \emph{Tests} names what the source puts to a test, which is not always the operation the later literature borrows. One of the eleven source works is not in our archive and carries no verdict here: \textbf{absence of a reading is not absence of validation}. It documents the clause carried on each band of the audit table in the main text.}
\label{tab:validation}
\scriptsize
\setlength{\tabcolsep}{2.6pt}
\renewcommand{\arraystretch}{0.96}
\begin{tabular}{@{}>{\raggedright\arraybackslash}p{2.1cm}>{\raggedright\arraybackslash}p{2.4cm}>{\raggedright\arraybackslash}p{3.5cm}>{\raggedright\arraybackslash}p{3.6cm}@{}}
\toprule
\rowcolor{headerblue}
\hd{Source} & \hd{Tests} & \hd{What it reports} & \hd{Verbatim excerpt} \\
\midrule
\rowcolor{perf5!35}\citet{ke2006highlevel} & \textbf{human judgement of the property} & community ratings on 60{,}000 photos as ground truth; per-feature precision--recall and a combined classifier at 27.8 per cent equal error, on the extremes only & \emph{``the photos have been rated by its community of users, giving us ground truth on the quality of each photo \,\ldots\, The top and bottom 10\% of the photos were extracted \,\ldots\, we chose to ignore the middle 80\%''} \\
\rowcolor{perf5!35}\citet{hasler2003colorfulness} & \textbf{human judgement of the property} & a psychophysical scaling experiment, seven categories, the expression fitted to the ratings at 94.0 per cent; the property is colourfulness & \emph{``We set up a psychophysical category scaling experiment, and ask people to rate images using 7 categories of colourfulness. We then fit a metric to the results''} \\
\rowcolor{perf2!45}\citet{srinivas2016rank} & the summary only & precision, recall and F of the selected keyframes against the ground-truth keyframes of seven open-video.org videos & \emph{``results are compared with the ground truth keyframes provided by open-video.org''} \\
\rowcolor{perf2!45}\citet{qaroush2025diversity} & the summary only & precision, recall and F1 of the extracted segments against reference summaries & \emph{``The evaluation focuses on comparing segments extracted by our system with those identified in ground truth summaries''} \\
\rowcolor{perf1!45}\citet{korolkov2025keyframe} & \textbf{neither the rule nor its weights} & weights set, not fitted; ablations vary segmentation parameters, and keyframes are judged by coverage and a qualitative review & \emph{``We set wsharp = 0.7, wbright = 0.3 to favor focused frames''} \\
\rowcolor{perf3!55}\citet{cakmak2025tripss} & the summary, with an ablation & TVSum20 and SumMe against human importance scores, with an ablation removing each modality & \emph{``Table 1 presents the results of an ablation study analyzing the contribution of each modality. Removing any single component''} \\
\bottomrule
\end{tabular}
\end{table}

%% GERADO por gera_tab_validacao.py. Nao editar a mao.
\begin{table}[p]
\centering
\caption{\textbf{Evidence reported by each source work, continued.} Columns as in the preceding table.}
\label{tab:validation_b}
\scriptsize
\setlength{\tabcolsep}{2.6pt}
\renewcommand{\arraystretch}{0.96}
\begin{tabular}{@{}>{\raggedright\arraybackslash}p{2.1cm}>{\raggedright\arraybackslash}p{2.4cm}>{\raggedright\arraybackslash}p{3.5cm}>{\raggedright\arraybackslash}p{3.6cm}@{}}
\toprule
\rowcolor{headerblue}
\hd{Source} & \hd{Tests} & \hd{What it reports} & \hd{Verbatim excerpt} \\
\midrule
\rowcolor{perf3!55}\citet{yang2024keyframe} & the summary, with an ablation & structural and feature ablations scored by F on SumMe over 15 annotations & \emph{``we conducted experiments on the SumMe dataset for structural ablation of the model and for feature ablation''} \\
\rowcolor{perf3!55}\citet{gygli2014summe} & the summary, with an ablation & f-measure against multiple human-created summaries per video, collected in a study with perception psychologists, plus an ablation measuring the gain of adding each feature to the set & \emph{``We evaluate our method using multiple `ground truth' summaries per video, which we acquired in a study in collaboration with perception psychologists''} \\
\rowcolor{perf3!55}\citet{jeong2016static} & the summary, with an optimisation ablation & summary-level human evaluation and precision/recall against labelled events, with a comparison against an optimisation-removed variant; the component-level claim here concerns the isolated grey-level statistic, which is not tested alone & \emph{``and (5) our method without the optimization procedure''} \textnormal{(Section 5, Experimental results)} \\
\rowcolor{perf2!45}\citet{huang2024aesthetic} & the summary only & Spearman and Kendall rank correlation against human annotations on four datasets, in canonical, augmentation and transfer setups & \emph{``we conduct an extensive evaluation of the models we developed by utilizing four distinct datasets: SumMe, TVSum, OVP, and YouTube \,\ldots\, we carry out our experiments using the dataset in three different setups: canonical (C), augmentation (A), and transfer (T)''} \textnormal{(Section IV.A, Experiments setup)} \\
\rowcolor{perf3!55}\citet{xu2024property} & the summary, with an ablation & F against summaries from five users on OVP, with an ablation of the frame-block length; no isolated test of the entities or of their weights & \emph{``All the scores are normalized to the interval [0, 1] and the higher value means the better''} \textnormal{(Section "Scoring the entities")} \\
\midrule
\bottomrule
\end{tabular}
\end{table}
